\documentclass[10pt,a4paper]{amsart}
\usepackage[margin=19mm]{geometry}
\usepackage{amsmath,amssymb,mathtools}
\usepackage[scr=rsfs,cal=euler]{mathalfa}
\usepackage{xfrac}
\usepackage[unicode,hidelinks,bookmarks=false]{hyperref}
\newcommand{\ie}{i.e.\ }
\newcommand{\cf}{cf.\ }
\newcommand{\resp}{respectively\ }
\newcommand{\Subsection}{\subsection}
\providecommand{\nobreakdash}{\nobreak\hbox{-}}
\setlength{\emergencystretch}{2em}
\multlinegap0pt
\usepackage{amssymb}
\usepackage{tikz}
\newtheorem{prop}{Proposition}
\newtheorem{lem}{Lemma}
\newcommand{\GL}{\mathrm{GL}}
\newcommand{\Gal}{\mathrm{Gal}}
\newcommand{\fp}{\mathfrak{p}}
\title{Frobenius Traces for Rank-2 Drinfeld Modules, Higher-Dimensional Galois Representations, and a Strong Multiplicity One Theorem in Positive Characteristic}
\author{Chien-Hua Chen}
\date{Source: arXiv:2604.28015v2 (CC BY 4.0)}
\begin{document}
\maketitle

\noindent\emph{Source and licence.} Frobenius Traces for Rank-2 Drinfeld Modules, Higher-Dimensional Galois Representations, and a Strong Multiplicity One Theorem in Positive Characteristic. Version arXiv:2604.28015v2 (CC BY 4.0), distributed by arXiv under the Creative Commons Attribution 4.0 International licence, http://creativecommons.org/licenses/by/4.0/, as recorded in the arXiv licence metadata for this exact version and shown on the abstract page https://arxiv.org/abs/2604.28015v2. Full licence notice: \texttt{licenses/arxiv-2604.28015-CC-BY-4.0.txt}.\par\smallskip

\noindent\textit{Editorial scope.} Counted unit: the article's prop Main1 (source0.tex lines 324--331). The article's complete original proof of it is delivered with it (source0.tex lines 333--409, one \texttt{proof} environment of 77 lines). No mathematics is written, repaired or re-derived: the statement and the proof are the article's own lines, in the article's own notation, and no retained line is substituted or re-flowed.  No further source-local context is needed: every cross-reference of every retained line resolves inside the two delivered spans.  Nothing was omitted from any retained span, so the package carries no omission record.  No retained line contains a citation, so the wrapper carries no bibliography and no bibliography entry.  No retained line contains a figure environment, an image-inclusion command or drawing code, so no image is delivered and the package declares no asset dependency.  Editorial formatting only, in addition: an AMS \texttt{amsart} preamble that restates the article's own declarations and macros used by the retained lines, namely the environments \texttt{prop}, \texttt{lem} on separate counters, together with the standard packages those lines need.  The retained environments print in this document as Proposition 1, Lemma 1 in order of appearance, whereas the article prints the same items with the numbering of its own full document.  The complete native source of the article as distributed in the arXiv e-print for this exact version is bundled unchanged as \texttt{sources/199518/source0.tex}.
\begin{prop}\label{Main1}
Let $K$ be a global function field, $G_K:=\Gal(\bar{K}/K)$, and $E$ be any local field of characteristic $p>0$. Let $n\in \mathbb{N}$, and $\rho_1, \rho_2:G_K\rightarrow \GL_n(E)$ be continuous representations, unramified outside a finite set $S$ of places of $K$. Suppose further that
\begin{itemize}
\item For any $ \fp\not\in S$, we have ${\rm Tr}(\rho_1({\rm Frob}_\fp))={\rm Tr}(\rho_2({\rm Frob}_\fp))$.
\item $\rho_1$ is absolutely irreducible and $\rho_2$ is irreducible.
\end{itemize}
Then $\rho_1\cong\rho_2$
\end{prop}

\begin{proof}

We start with the fact that for $i=1,2$,  the composition $${\rm Tr}\circ \rho_i: G_K\rightarrow E\ \ \ g\mapsto {\rm Tr}(\rho_i(g))$$ is
continuous and invariant under conjugation. Hence Chebotarev density Theorem implies that ${\rm Tr}(\rho_1(g))={\rm Tr}(\rho_2(g))$ for any $g\in G_K$.




Now we consider the well-known Burnside's theorem below:
\begin{lem}[Burnside's Theorem]
Let $E$ be a field, $V$ be a finite dimensional vector space over $E$, $K$ be a  field with $G_K:=\Gal(\bar{K}/K)$, and $\rho:G_K\rightarrow GL_E(V)$ be a representation. Set $A:=E[\rho(G_K)]\subset {\rm End}_E(V)$. If $\rho$ is absolutely irreducible, then $A={\rm End}_E(V)$.
\end{lem}

Set $A_1=E[\rho_1(G_K)]={\rm End}_E(V_1)$ and $A_2=E[\rho_2(G_K)]$, where $V_1$ is a $n$-dimensional vector space over $E$ with $G_K$-action defined via $\rho_1$, similar for the vector space $V_2$. We fix a $E$-basis $\{a_i\}_{i=1}^{n^2}$ of $A$, where $a_i=\rho_1(g_i)$ for some $g_i\in G_K$. Let $\{a_i^*\}_{i=1}^{n^2}$ be the trace-dual basis of ${\rm End}_E(V_1)$, i.e. ${\rm Tr}(a_i^*a_j)=\delta_{ij}$. Then we define the $E$-linear map

$$
\begin{array}{lrl}
\Pi:&{\rm Hom}_E(V_1,V_2)&\rightarrow {\rm Hom}_E(V_1,V_2)\\
&S& \mapsto \sum_{i=1}^{n^2}\rho_2(g_i)Sa_i^*
\end{array}
$$
On the other hand, we define the following $E$-algebra homomorphism

$$\Phi:{\rm End}_E(V_1)\rightarrow E[\rho_2(G_K)]\ \ \sum_{i=1}^{n^2} c_ia_i\mapsto \sum_{i=1}^{n^2}c_i\rho_2(g_i).$$

Note that $\Phi$ is well-defined since for $\sum_{i=1}^{n^2}c_ia_i=0$, we have  ${\rm Tr}\left(\left [\sum_{i=1}^{n^2}c_i \rho_1(g_i)\right ]\rho_1(g)\right)=0$ for any $g\in G_K$. Its image via $\Phi$ has trace 
$${\rm Tr}\left(\left [\sum_{i=1}^{n^2}c_i \rho_2(g_i)\right ]\rho_2(g)\right)=0 \textrm{ for all $g\in G_K$}.$$ 
As $\rho_2$ is semisimple, $A_2$ is a semisimple algebra, thus the trace pairing on $A_2$ is non-degenerate, therefore the above trace equality implies $\sum_{i=1}^{n^2}c_i \rho_2(g_i)=0$. 



Now we claim that for $S\in {\rm Hom}_E(V_1,V_2)$, $\Pi(S):V_1\rightarrow V_2$ is $G_K$-equivariant, and $\Pi$ is not a zero map. Given a fixed $g\in G_K$, we can write 
\begin{equation}
\rho_1(g) a_i = \sum_{j=1}^{n^2} \text{Tr}(a_j^* \rho_1(g) a_i) a_j \tag{1}, \textrm{ and set }c_{ij}(g)=\text{Tr}(a_j^* \rho_1(g) a_i).\end{equation}

Similarly, we have \begin{equation}
a_j^* \rho_1(g) = \sum_{i=1}^{n^2} \text{Tr}(a_j^* \rho_1(g) a_i) a_i^* \tag{2}, \textrm{ and set }d_{ji}=\text{Tr}(a_j^* \rho_1(g) a_i).\end{equation}
One can easily see that $c_{ij}(g)=d_{ji}(g)$. Now we compute
$$
\begin{array}{ll}
\rho_2(g)\Pi(S)&=\sum_{i=1}^{n^2}\rho_2(gg_i)Sa_i^*=\sum_{i=1}^{n^2}\Phi(\rho_1(g)a_i)Sa_i^*=\sum_{i=1}^{n^2}\sum_{j=1}^{n^2}c_{ij}(g)\Phi(a_j)Sa_i^*\\
&\ \\
&=\sum_{i=1}^{n^2}\sum_{j=1}^{n^2}c_{ij}(g)\rho_2(g_j)Sa_i^*=\sum_{j=1}^{n^2}\rho_2(g_j)S\left(\sum_{i=1}^{n^2} d_{ji}(g)a_i^*\right)\\
&\ \\
&=\sum_{j=1}^{n^2}\rho_2(g_j)Sa_j^*\rho_1(g)=\Pi(S)\rho_1(g).
\end{array}
$$
It remains to prove $\Pi$ is not a zero map, we start with the operator $\Pi(S)=\sum_i \rho_2(g_i)\, S\, a_i^\ast.$
Consider the equality $\Phi(a_i)=\rho_2(g_i)$. We have seen that $\Phi:{\rm End}_E(V_1)\rightarrow E[\rho_2(G_K)]\subset {\rm End}_E(V_2)$ is an $E$-algebra homomorphism. As ${\rm End}_E(V_1)$ is a simple algebra, $\Phi$ must be injective. Comparing $E$-dimension on ${\rm End}_E(V_i)$ for $i=1,2$ forces $\Phi$ to be an $E$-algebra isomorphism. Thus apply Skolem--Noether theorem we may write
\[
\Phi(a_i)=Ta_iT^{-1}
\]
for some $E$-linear isomorphism $T:V_1 \to V_2$. Hence we obtain
\[
\Pi(S)=\sum_i T a_i T^{-1} S a_i^\ast.
\]
Now choose $S = TX$ for some $X \in {\rm End}_E(V_1)$. Then
\[
\Pi(TX)
=
\sum_i T a_i T^{-1} T X a_i^\ast
=
T\left(\sum_i a_i X a_i^\ast\right).
\]
Using the identity
$$\sum_i a_i X a_i^\ast = {\rm Tr}(X)\,\mathrm{Id},$$
we get $\Pi(TX)={\rm Tr}(X)\,T.$ Put $X=E_{11}$ to be the elementary matrix, we have ${\rm Tr}(E_{11})=1$, and hence
\[
\Pi(T E_{11}) = T \neq 0.
\]
This proves the claim.

Finally, we choose a  non-zero map $S$ such that $\Pi(S)\neq0$, then apply Schur's lemma on the nonzero $G_K$-module homomorphism $\Pi(S):V_1\rightarrow V_2$ to get an isomorphism between $V_1$ and $V_2$, hence $\rho_1\cong\rho_2$.



\end{proof}

\end{document}
